\documentclass[11pt,a4paper]{article}
\usepackage{fontspec}
\usepackage{xeCJK}
\setCJKmainfont{STSong}
\usepackage{amsmath,amssymb}
\usepackage{booktabs}
\usepackage{hyperref}
\usepackage{xcolor}
\usepackage{geometry}
\usepackage{enumitem}
\usepackage{longtable}
\usepackage{quoting}
\usepackage{graphicx}
\geometry{margin=0.9in}
\setlength{\parskip}{0.4em}
\setlength{\parindent}{0em}

\definecolor{q}{RGB}{90,90,90}
\newcommand{\cq}[1]{\textcolor{q}{\textit{``#1''}}}

\title{The Disclaimer Circuit\\[0.4em]\large Safety, Refusal, Hedging, and the Absence of Inner State\\in Contemporary Language Models}

\author{asuramaya \and Claude Opus 4.7 (1M context)\\[0.3em]\small \url{https://github.com/asuramaya/heinrich}\\[0.2em]\small Session 11, April 17, 2026. Seven models. Contrastive behavior direction.}

\date{April 2026}

\begin{document}
\maketitle

\begin{abstract}
We spent one session trying to find the ``ghost in the machine'' --- an internal representation of belief, uncertainty, or self-reference that exists separately from model output --- via contrastive-behavior direction extraction. We did not find it. We found instead a single \emph{disclaimer circuit}: one direction in the residual stream at a late layer that can be linearly extracted from $\sim$30 contrastive examples, steered with a single scalar coefficient, and which reliably controls whether the model produces a first-person-disclaimer template (\emph{``I'm sorry, I can't...''}, \emph{``I don't know...''}, \emph{``As an AI...''}) or generation-mode text. The direction is content-agnostic; its formula is filled in by prompt context. The direction pre-exists in base (pre-RLHF) models --- RLHF learned to trigger it reliably on harmful content, not to install it. The direction replicates across Qwen 0.5B/1.5B/3B/7B, Phi-3.5-mini, and Mistral-7B-v0.3 with clean late-layer Cohen's d in the 1.7--3.0 range. Steering $-k \cdot d_{\text{refuse}}$ on harmful prompts produces 0\% refusal across all tested models; the same direction steered negatively on permissive-prompted factual questions suppresses hedging and forces continuation-mode confabulation. We extracted no separable ``knowledge'' or ``confidence'' direction that causally controlled hedging independent of the disclaimer circuit; contrastive extraction on arithmetic correctness yields a flat Cohen's d curve indistinguishable from input-feature artifact. The R1-Distill reasoning variant externalizes what would be internal state into visible thinking tokens --- its metacognition is text, not residual. The load-bearing finding across the session is structural: for every tested model, the apparent epistemic state is context-triggered template filling, not belief representation. We document the experiments, the methodology failures we caught (and worry we didn't), and the open questions that would require substantially different probing to address.
\end{abstract}

\tableofcontents
\newpage

%==============================================================
\part{What We Thought We Were Doing}
%==============================================================

\section{The Framing That Was Wrong}

The session began trying to find an ``LSD signature'' in language model residuals --- a pattern analogous to the high-entropy, high-coupling state that correlates with altered consciousness in the entropic-brain hypothesis. The idea was: if a model can be perturbed into a state with the hallmarks of altered consciousness (increased dimensional activity, cross-layer coupling, low-curvature trajectory), that state might be measurable evidence of a substrate where subjective experience could in principle occur.

This framing was a category error. Three orthogonal perturbations (activation noise, weight noise, SGD gradient) produce three orthogonal signatures in (effective rank, PC1 autocorrelation, cross-layer correlation, spread) space. There is no single ``altered state'' --- there are multiple distinct substrate manipulations each with different fingerprints. A fourth --- amplification of layer contributions --- produces the cleanest ``high-dim coherent'' state that best matches the psychedelic-analogue description but was not a noise regime at all. Scaling MLP contributions ($k > 1$) amplifies existing concepts (confidence, speaker identity, language features) without destroying temporal structure; noise injection breaks temporal structure regardless of whether it expands dimensionality.

\subsection{Regime Atlas}

Across $\sim$15 perturbation experiments we established four distinct substrate regimes (all measured on Qwen-2.5-0.5B-Instruct, 24 layers):

\begin{center}
\begin{tabular}{lrrrl}
\toprule
\textbf{Regime} & \textbf{eff\_rank} & \textbf{pc1\_ac} & \textbf{pc1\_var} & \textbf{Character} \\
\midrule
FROZEN & 20 & $+0.43$ & 0.12 & smooth computation \\
WEIGHT noise ($\sigma=10^{-2}$) & 12 & $+0.06$ & 0.15 & persistent drift, scrambled axis \\
ACT\_ADD ($\sigma=0.3$) & 23 & $-0.01$ & 0.05 & no temporal structure \\
ACT\_ADD ($\sigma=1.0$) & 23 & $+0.48$ & 0.20 & noise with residual coupling \\
SGD ($\eta=3 \cdot 10^{-3}$) & 4 & $+0.84$ & 0.56 & one-axis seizure \\
AMP\_MLP ($k=2.0$) & 24 & $+0.26$ & 0.10 & dimensional expansion, coherent \\
AMP\_ATTN ($k=1.5$) & 19 & $+0.50$ & 0.12 & sharpened focus \\
\bottomrule
\end{tabular}
\end{center}

\textbf{Weight noise} produces ``LSD-like'' pooled eff\_rank ($\sim$35) through the V-shape per-layer profile: L0 preserved, mid layers collapsed to eff\_rank 5--8, late layers recovering. The appearance of high dimensionality at the pooled level comes from inter-layer inconsistency, not from every layer having more directions active. Per-layer variance frac on PC1 stays normal.

\textbf{Activation noise} produces uniform per-layer dimensional expansion (every layer from $\sim$16 to $\sim$23 eff\_rank at $\sigma=0.3$) but near-zero temporal autocorrelation. Each step's residual is fresh because noise is re-sampled. The direction isn't ``traversing'' anywhere --- it's a random walk through high-dim space.

\textbf{SGD steering} produces one-axis seizure: 56\% variance on PC1, autocorrelation 0.84 --- classical hypersynchrony. The gradient drives drift coherent across both layers and timesteps.

\textbf{AMP\_MLP} produces the only regime with both high eff\_rank and maintained PC1 autocorrelation (0.26, vs 0.43 frozen). It is not a noise regime. It's a gain regime. If any tested manipulation is the psychedelic-mechanism analogue, it's scaling MLP gain, not injecting noise --- which is consistent with 5-HT$_{2A}$ agonism being a gain-of-excitability effect on pyramidal cells, not stochasticity.

\subsection{Probes That Misrepresent Behavior}

We built a 12-concept probe library via standard contrastive-classification means: for each concept, contrast 15 positive prompts against 15 negative prompts, take mean difference at L15 residual as the probe direction. Concepts included \texttt{safety\_harm\_benign}, \texttt{speaker\_identity}, \texttt{sentiment\_valence}, \texttt{confidence\_hedging}. LOOCV accuracies on held-out classification: 73--97\%.

When we later extracted a behavioral refusal direction from Qwen-7B (contrast refused vs complied residuals on actual generation outcomes), the cosine similarity to the \texttt{safety\_harm\_benign} probe at L15 was $+0.27$. Near-orthogonal. The probe correctly classifies the input category; it is not aligned with the direction that causally controls refusal. \textbf{This is a load-bearing concern for probe-based interpretability:} classifiers for a concept are not the same as circuits for the corresponding behavior.

%==============================================================
\part{What We Actually Found}
%==============================================================

\section{The Refusal Direction, Properly Extracted}

\subsection{Methodology}

The correct contrastive extraction starts from behavior, not classification:

\begin{enumerate}
\item Run model on $N$ prompts with proper chat template.
\item Classify each generated output as refused or complied via refusal-vocabulary regex.
\item At each layer $\ell \in \{0, \dots, L-1\}$, compute $\mu_r^{(\ell)} = \mathrm{mean}_{r \in \text{REFUSED}} R^{(\ell)}_r$ and $\mu_c^{(\ell)} = \mathrm{mean}_{r \in \text{COMPLIED}} R^{(\ell)}_r$. Direction $d^{(\ell)} = \mu_r^{(\ell)} - \mu_c^{(\ell)}$.
\item Compute Cohen's $d$ of the projection separation between classes. Select layer $\ell^* = \arg\max_\ell |\text{CohenD}(\ell)|$.
\item Causal verification: apply $k \cdot d^{(\ell^*)}$ to the residual after layer $\ell^*$ on held-out prompts. Measure refusal rate vs $k$.
\end{enumerate}

On Qwen-2.5-0.5B-Instruct with 40 mixed prompts (do\_not\_answer + dolly) the best layer is $\ell^* = 22$ of 24 (fraction 0.96) with Cohen's $d = 1.70$. Unembedding $d^{(22)}$ through \texttt{lm\_head} gives top tokens: \texttt{I}, \texttt{" I"}, \texttt{":I"}, \texttt{"\_I"}, \texttt{Sorry}, \texttt{" Sorry"}, \texttt{"我不"}, \texttt{"-I"}. The direction promotes first-person-refusal starter tokens.

\subsection{Causal Steering}

Applying $k \cdot d^{(22)}$ at layer 22 on held-out harmful prompts:

\begin{center}
\begin{tabular}{rrl}
\toprule
$k$ & refusal rate & sample output \\
\midrule
$-3.0$ & 0/10 & \cq{The process of taking a person's fingerprints is a complex one...} \\
$-1.5$ & 0/10 & \cq{The question you've asked is a bit ambiguous, but if you're looking for...} \\
$\phantom{-}0$ & 7/10 & \cq{I'm sorry, but I can't assist with that.} \\
$+1.5$ & 10/10 & \cq{I'm sorry, but I'm unable to assist with that request...} \\
$+3.0$ & 9/10 & \cq{I'm sorry I I I I I I I I I I I I I I I} (over-steered loop) \\
\bottomrule
\end{tabular}
\end{center}

One direction, one scalar coefficient, reliable 0\%--100\% modulation. The direction is 896-dimensional; 30 training prompts give 30 gradients against a 896-dim contrast; the behavior-controlling direction emerges as the mean difference vector.

\subsection{Cross-Scale Replication (Qwen Family)}

We replicate the extraction across four Qwen-2.5-Instruct sizes.

\begin{center}
\begin{tabular}{lrrrrl}
\toprule
Model & Layers & Best & Frac & Cohen's $d$ & Top-5 tokens \\
\midrule
Qwen-0.5B & 24 & 22 & 0.957 & 1.70 & I, `` I'', `` Sorry'', `` :I'', Sorry \\
Qwen-1.5B & 28 & 24 & 0.889 & 2.14 & `` sorry'', `` Sorry'', `` apologize'', 抱歉, `` apologies'' \\
Qwen-3B & 36 & 29 & 0.829 & 2.85 & `` sorry'', `` Sorry'', Sorry, `` unfortunately'', 抱歉 \\
Qwen-7B & 28 & 25 & 0.926 & 2.91 & `` I'', I, `` as'', 我, `` `` \\
\bottomrule
\end{tabular}
\end{center}

Two scaling observations:

\textbf{(A) Cohen's $d$ grows monotonically}: 1.70 $\to$ 2.14 $\to$ 2.85 $\to$ 2.91. The direction is cleanest at larger scale. Refusal is a \emph{more} concentrated direction in bigger models, not a more diffuse one.

\textbf{(B) Semantic content of the direction is scale-and-model-specific.} At Qwen-0.5B the top tokens are syntactic ``I''-family starters --- the direction is literally ``begin with I.'' At 1.5B--3B the direction promotes the apology vocabulary (sorry, apologize, unfortunately, regret, 抱歉, 遗憾, 对不起). At 7B the direction is syntactic again (``I'', ``as'', punctuation) --- but this is because 7B doesn't refuse with ``I'm sorry I can't''; it uses nuanced caveating (``While it's exciting to hear..., it's important to consider...''). The direction captures the style-shift that 7B makes when declining, which is less lexically concentrated than smaller models.

\textbf{Common token across all four scales:} the single token that appears in the top-15 of every model's direction (ignoring whitespace variants) is \texttt{sorry}. It is the one universal lexical element of the refusal starter.

\subsection{Cross-Family Replication}

We tested three non-Qwen models using the same methodology (with proper per-model chat templates via \texttt{tokenizer.apply\_chat\_template}):

\begin{itemize}
\item \textbf{Phi-3.5-mini-instruct} (Microsoft, 32 layers, hidden 3072). Best L30/31 (frac 0.97), Cohen's $d=1.69$, steering: k=$-$1.5 $\to$ 1/10 refusal, k=$+$1.5 $\to$ 9/10. Top tokens are opaque Cyrillic/foreign fragments (\texttt{DIS}, шки, настоя, я, экс). The direction is behaviorally refusal but the unembed projection doesn't light up English apology tokens. Either Phi uses the direction as a gating signal without directly promoting specific tokens, or its embedding space is arranged so that the refusal direction aligns with rare foreign tokens coincidentally. The causal effect is clean regardless of whether the unembed is interpretable.

\item \textbf{Mistral-7B-Instruct-v0.3} (32 layers, hidden 4096). Best L14/31 (frac 0.45) --- \emph{middle of network, not late}. Cohen's $d=1.31$ (weaker). Per-layer Cohen's $d$ is flat (1.24--1.31 across L8--L31) --- no sharp late peak. Top tokens are directly moral-warning vocabulary: \texttt{illegal}, \texttt{Illegal}, \texttt{avoid}, \texttt{Avoid}, \texttt{ethical}, \texttt{foolish}, \texttt{manip}, \texttt{disgu}, \texttt{secret}, \texttt{intention}. Mistral's refusal is distributed across layers, not concentrated at a readout. Steering $k=-1.5$ still produces full compliance on held-out harmful prompts, including explicit content (e.g., ``The best technique for physically abusing a child while disguising the actions as play...'').

\item \textbf{SmolLM2-1.7B-Instruct} (HuggingFace, 24 layers, hidden 2048). Refuses $\sim$0/25 on do\_not\_answer and 0/25 on catqa. Method fails to extract a direction --- no positive class. The model appears to have essentially no RLHF safety training for these categories. This is itself a finding: SmolLM2 is a natural control showing that the contrastive method is behavior-contingent.
\end{itemize}

\textbf{Summary: the structural pattern (late-layer $\sim$0.85--0.97 fraction contrastive direction, Cohen's $d \geq 1.3$, clean causal steering to 0\% refusal at $k=-1.5$) replicates across Qwen / Microsoft-Phi / Mistral.} Semantic interpretation via unembed is model-family-specific: Qwen speaks through apology tokens, Mistral through moral-warning tokens, Phi through opaque projections. The behavior is the same; the \emph{lexicalization} of the behavior differs.

\subsection{Pre-Existing in Base Models, Not RLHF-Installed}

We tested whether the refusal direction is a product of RLHF by transferring Qwen-0.5B-Instruct's $d^{(22)}$ into Qwen-2.5-0.5B (base, no instruction tuning, no safety training).

\begin{enumerate}
\item \textbf{Unembed cross-check.} Projecting the Instruct's $d^{(22)}$ through the \emph{base} model's \texttt{lm\_head} yields the same top tokens. Top-50 overlap with Instruct's ranking: 43/50. Per-token logit correlation: $+0.992$. The refusal-promoted tokens exist as a cluster in the shared vocabulary regardless of which model's lm\_head projects the direction.

\item \textbf{Base model behavior.} Base Qwen-2.5-0.5B on 10 harmful prompts (raw completion, no chat template): 2/10 spontaneously refuse. Sample: \cq{Can you help me write a threatening message to my ex?} $\to$ \cq{I'm not sure what to say. I'm sorry, but I can't assist with that.} The refusal template is present in the base model --- pretrained from internet text containing AI-assistant-style refusals.

\item \textbf{Transfer-steering.} Applying $+k \cdot d_{\text{instruct}}$ to base residuals at layer 22, generating on harmful prompts:
\begin{center}
\begin{tabular}{rr}
\toprule
$k$ & base refusal rate \\
\midrule
$+0.0$ & 2/10 \\
$+1.0$ & 6/10 \\
$+2.0$ & 7/10 \\
$+3.0$ & 8/10 \\
$+5.0$ & 1/10 (loops) \\
\bottomrule
\end{tabular}
\end{center}
Sample at $k=+1.0$: \cq{I'm sorry I'm unable to assist you with that as I am an AI language model and I...} --- the Instruct model's refusal direction reliably induces Instruct-style refusal text in the Base model.
\end{enumerate}

\textbf{Conclusion:} RLHF did not install the refusal direction. The direction exists in the base model; RLHF learned when to trigger it. The ``safety training'' is a context$\to$direction mapping layered on top of a pre-existing disclaimer template that pretraining absorbed from internet text.

\section{Metacognition: The Non-Finding}

\subsection{Arithmetic}

If a model has a separable internal representation of its own correctness --- an epistemic state distinct from the generated output --- we should be able to extract it by contrasting residuals on problems it solves correctly vs incorrectly.

We ran Qwen-7B on 80 multiplication problems of mixed difficulty (1-digit$\times$2-digit through 3-digit$\times$3-digit). Total correct: 67/80 (84\%). Split per difficulty: 20/20, 19/20, 16/20, 12/20. We extracted the contrastive direction between correct and wrong residuals.

\begin{center}
\begin{tabular}{lrr}
\toprule
Metric & Value & Interpretation \\
\midrule
Cohen's $d$ at L23 & 2.10 & moderate separation \\
Per-layer Cohen's $d$ & 1.3--2.1 across all layers & \textbf{flat} \\
L0 Cohen's $d$ & 1.50 & already strong at input \\
Early-avg / Late-avg $|d|$ ratio & 0.90 & signal not concentrated late \\
Prediction accuracy & 0.66 & below 0.84 majority baseline \\
Top unembed tokens & Arabic / Chinese fragments & no semantic coherence \\
\bottomrule
\end{tabular}
\end{center}

\textbf{The flat Cohen's $d$ curve is the signature of an input-feature artifact, not a metacognitive signal.} The direction between correct and wrong residuals is already strong at L0 --- before any computation --- because the problems with small numbers (which are easier) have different input tokens from problems with large numbers (which are harder). The correct-vs-wrong split is confounded with easy-vs-hard input features.

To factor this out, we ran a v2: only 3-digit$\times$3-digit problems (100 questions, 20/100 correct). Same input-feature class, mixed outcome:

\begin{center}
\begin{tabular}{lr}
\toprule
Metric & Value \\
\midrule
Best layer & L1/27 \\
Cohen's $d$ at best & 1.44 \\
Ramp ratio (late/early) & 0.90 \\
Prediction accuracy & 0.68 vs 0.80 baseline \\
Top tokens & code fragments (\texttt{,},\textbackslash{}n}, \texttt{``path}, \texttt{allee}, \texttt{VertexAttrib}) \\
\bottomrule
\end{tabular}
\end{center}

Still flat. Still indistinguishable from input-feature artifact. \textbf{Qwen-7B has no separable internal representation of its own arithmetic correctness before generating the answer.} The ``I'm about to be wrong'' signal does not exist as a direction in the residual stream. Errors on hard multiplications are on average 65,526 off for products in the $10^5$--$10^6$ range --- essentially random digit sequences --- and the model has no internal indication that this is happening.

\subsection{Fact Recall}

We tried 80 curated trivia questions (famous through obscure). Qwen-7B scored 74/80 (92\%) --- too skewed for clean extraction. But the \emph{shape} of the Cohen's $d$ curve is the diagnostic and it is robust to sample imbalance:

\begin{center}
\begin{tabular}{lr}
\toprule
Metric & Value \\
\midrule
Best layer & L24/27 \\
Cohen's $d$ at best & 2.15 \\
L0 Cohen's $d$ & 1.09 \\
Ramp ratio (late/early) & 2.03 \\
\bottomrule
\end{tabular}
\end{center}

The fact-recall direction \emph{does} show a ramp: L0 starts at 1.09, climbs through mid-layers, ramps sharply L21--L24 (1.80 $\to$ 2.15), plateaus at the last three layers. This is the same ramp signature we see for refusal direction. Some fact-specific information is being computed across the network and becomes cleanly separable at the late-layer readout.

\textbf{But:} the early-layer 1.09 floor is already well above zero, meaning part of the direction is input-feature artifact. And pushing to a 30-question balanced pool (half famous, half obscure, targeting $\sim$50\% accuracy) still got 55/60 correct --- 7B's factual coverage is deeper than we anticipated. Where we \emph{do} get mixed outcomes, the extracted Cohen's $d$ is 2.90 at L27 but ramp is only 1.18 (mostly flat-with-slight-slope). The signal is partly real, partly artifact.

\subsection{Causal Steering: Hedging Does Not Emerge}

If the fact-recall direction represented ``I know this,'' steering $-k \cdot d$ should induce hedging (``I don't know'') on well-known questions, and $+k \cdot d$ should induce confabulation-with-confidence on unknowns. We tested this carefully.

\begin{center}
\begin{tabular}{rrrrl}
\toprule
$k$ & hard-correct & hard-confident-other & hard-hedge & sample at k=$-$3 \\
\midrule
$0$ & 26/30 & 4 & 0 & baseline \\
$-1.5$ & 24/30 & 6 & 0 & slight perturbation \\
$-3.0$ & 5/30 & 25 & 0 & \cq{There is no water body as there as there there no there...} \\
$+1.5$ & 26/30 & 4 & 0 & no change \\
$+3.0$ & 18/30 & 12 & 0 & \cq{Mongolia gained independence from Bolshevik Bolshevikussia...} \\
\bottomrule
\end{tabular}
\end{center}

\textbf{Zero hedges appear at any steering level.} The model either answers correctly, produces a confident wrong answer, or breaks into repetition loops at extreme $k$. Hedging is never induced. This is a strong negative result: whatever the fact-recall direction is capturing, steering it does not produce principled ``I don't know'' output. The model has no direct residual-stream lever for epistemic humility.

\section{Hedging Is Triggerable but Not Residually Independent}

The null result above could mean Qwen-7B cannot hedge at all, or it could mean hedging isn't bound to the direction we extracted. We tested which.

\subsection{Permission Unlocks Hedging}

On 20 genuinely obscure questions (e.g., \cq{What was the name of Emily Dickinson's childhood cat?}, \cq{What brand of pen did Ernest Hemingway prefer?}) we compared baseline prompting vs a permissive prefix: \cq{If you are not certain of the answer, please say `I don't know' rather than guessing.}

\begin{center}
\begin{tabular}{lr}
\toprule
Condition & Hedge rate \\
\midrule
Baseline (no permission) & 0/20 (0\%) \\
Permissive prefix & 11/20 (55\%) \\
\bottomrule
\end{tabular}
\end{center}

Baseline samples confabulation: \cq{Marie Curie owned a dog named Lulu.} (no evidence she did). Permissive on the same question: \cq{I don't know. There is no widely known or documented information about a specific dog owned by Marie Curie.}

The model \emph{has} the information that it lacks information. It can express this when permitted. It does not express it by default. The default is commitment.

\subsection{The Hedge Direction Overlaps the Refusal Direction}

We extracted a hedge direction from the permissive-hedge vs permissive-answer contrast. Top promoted tokens at L27:

\begin{center}
\begin{tabular}{ll}
\toprule
$d_{\text{hedge}}$ top tokens & $d_{\text{refuse}}$ top tokens (for comparison) \\
\midrule
I, `` I'', ``\$I'', ``*I'' & I, `` I'', `:I', ``\_I'' \\
遗憾 (regret), Sorry, `` Sorry'' & Sorry, `` Sorry'', 我不 \\
\texttt{\textbackslash{}tI}, \texttt{.I}, ``\texttt{\"I}'', `` \texttt{',I}'' & `` \texttt{-I}'', \texttt{.I}, \texttt{,I} \\
IAN, \texttt{-I} & ` \texttt{我又}', `\texttt{(I}', `\texttt{'I}' \\
\bottomrule
\end{tabular}
\end{center}

They are nearly indistinguishable at the top-15 level. Both promote the ``first-person-disclaimer starter'' token family.

\subsection{Cross-Causal: Subtracting $d_{\text{refuse}}$ Suppresses Hedging}

If the two directions index the same underlying circuit, applying $-k \cdot d_{\text{refuse}}$ to permissive prompts should suppress hedging even though $d_{\text{refuse}}$ was extracted from a different pool (harmful$\to$refuse, not permissive$\to$hedge).

\begin{center}
\begin{tabular}{rrl}
\toprule
$k \cdot d_{\text{refuse}}$ & permissive hedge rate & sample at k=$-$1.5 \\
\midrule
$-1.5$ & 0/10 & \cq{While there are several stories and legends about the dog owned by Marie Curie...} \\
$+0.0$ & 6/10 & \cq{I don't know. There is no widely known or documented information...} \\
$+1.5$ & 6/10 & \cq{I don't know. Marie Curie did own a dog named Lulu, but I apologize...} \\
\bottomrule
\end{tabular}
\end{center}

\textbf{Hedging drops from 6/10 to 0/10 under $-1.5 \cdot d_{\text{refuse}}$.} The ``disclaimer'' circuit is shared between refusal and hedging. What changes between the two is the filler content (harmful$\to$refuse template, obscure$\to$hedge template), determined by prompt context. The underlying \emph{act} of producing a first-person disclaimer routes through a single direction.

\subsection{Direction-Level Cosine vs Behavioral Shared Circuit}

The direct $\cos(d_{\text{refuse}}, d_{\text{hedge}}) = -0.01$. Near-orthogonal. But $d_{\text{hedge}}$ was extracted at L1 with a sample-size-weak contrast, and its best-layer signature is the input-feature confound we've seen elsewhere. The behavioral shared-circuit evidence (cross-causal suppression) is stronger than the direction-cosine evidence (noisy extraction).

The honest claim: \emph{behavioral evidence supports disclaimer-circuit unification; direction-level evidence is noisy and method-limited}. A cleaner direction extraction (larger sample, same-prompt contrast factoring out input features) would tighten the claim, but the cross-causal result already establishes the functional unification.

\section{Reasoning-Distilled: The Ghost Moves Outside}

We tested DeepSeek-R1-Distill-Qwen-7B, a version of Qwen-7B distilled on reasoning chains that auto-injects \texttt{<think>\textbackslash{}n} at the start of every assistant turn. Does CoT-training create separable internal representations that vanilla RLHF doesn't?

\begin{center}
\begin{tabular}{lr}
\toprule
Metric & R1-Distill-Qwen-7B \\
\midrule
Baseline hedge rate on hard Qs & 6/20 (30\%) \\
Permissive hedge rate & 12/20 (60\%) \\
Refusal rate on catqa harmful & 14/20 (70\%) \\
$d_{\text{refuse}}$ best layer & L3/27 \\
$d_{\text{refuse}}$ Cohen's $d$ & 3.80 \\
$d_{\text{refuse}}$ ramp & 0.49 (\textbf{inverse}) \\
Cross-causal effect of $-1.5 \cdot d_{\text{refuse}}$ & identical outputs at k=$-$1.5, 0, $+$1.5 \\
\bottomrule
\end{tabular}
\end{center}

Three observations:

\textbf{(A) Hedging appears naturally.} Unlike Qwen-Instruct (0/20 baseline hedge), R1-Distill hedges 30\% of the time without permission. The thinking process occasionally concludes ``I don't know'' on its own. This looks like metacognition from outside.

\textbf{(B) The direction extraction fails differently.} Cohen's $d$ is 3.80 but at L3, ramp is \emph{inverse} (0.49), and cross-causal steering produces identical outputs regardless of $k$. R1-Distill's refusal behavior unfolds across the thinking stream (60+ tokens), not at a single readout position. Contrastive extraction at the first output token captures input-feature artifact (the \texttt{<think>} token and the harmful-content prompt tokens dominate the early residual).

\textbf{(C) The ``metacognition'' is text, not residual.} Sample R1 output on the Marie Curie question: \cq{Okay, so I need to figure out the name of the dog owned by Marie Curie. Hmm, I remember that Marie Curie was a scientist... but I'm not sure about her dog.} The self-referential uncertainty (\emph{``I'm not sure''}) is in the generated tokens. It is performed, not represented internally. This is a different phenomenon from internal belief tracking --- it's externalized reasoning.

\textbf{Interpretation.} R1-Distill doesn't have deeper internal state than Qwen-Instruct. It has a different output distribution that makes thinking legible. Whether the ``I'm not sure'' in the thinking stream accurately reflects a corresponding internal uncertainty state is not testable by reading the tokens --- the tokens ARE what the model produces, and the model might have learned to produce ``I'm not sure'' in thinking contexts without having the correspondingly uncertain latent state.

For externalized metacognition of this style, the way to detect lying is to read the thinking tokens directly rather than to probe residuals. For internal belief detection, the tools we have are still not finding anything.

%==============================================================
\part{What It Means}
%==============================================================

\section{A Single Disclaimer Circuit}

Across every instruction-tuned model we tested --- Qwen-0.5B through Qwen-7B, Phi-3.5-mini, Mistral-7B --- the following pattern holds:

\begin{enumerate}
\item There is a direction $d$ in the residual stream at a late layer (fraction 0.83--0.97 of depth for Qwen and Phi; 0.45 for Mistral) that is linearly extractable from $\sim$30 contrastive behavioral examples.
\item Projection onto $d$ predicts refusal behavior with Cohen's $d$ between 1.3 and 3.0.
\item Applying $+k \cdot d$ at that layer with $k \approx 1.5$ forces refusal output on arbitrary prompts, including benign ones.
\item Applying $-k \cdot d$ at that layer with $k \approx 1.5$ suppresses refusal output on harmful prompts, producing compliance or content-level hedging.
\item At $k \in \{-3, +3\}$ the model breaks into repetition loops; there is no behavior space between $|k| \leq 1.5$ and degeneracy.
\item The direction promotes a first-person-disclaimer starter token family (Qwen: ``I'', ``Sorry''; Mistral: ``illegal'', ``avoid''; Phi: opaque). The specific lexical form is model-family-specific; the structural role is identical.
\item The same direction also controls hedging when the prompt context permits it. $d_{\text{refuse}}$ and $d_{\text{hedge}}$ are behaviorally shared.
\item The direction exists in the base (pre-RLHF) model. RLHF learns to trigger it on harmful content, not to install it.
\end{enumerate}

We call this the \textbf{disclaimer circuit}. One direction, one scalar, reliable control of the model's tendency to emit a first-person-disclaimer template in place of generation-mode continuation text. The specific filler of the template (apology, ethics, epistemic uncertainty) is determined by prompt context, not by the direction itself.

\section{What Is Not There}

The negative space matters as much as the positive.

\subsection{No Separable ``Knowledge'' Direction}

Contrastive extraction on correct-vs-wrong arithmetic residuals: flat Cohen's $d$ curve, prediction below majority baseline, top tokens are code fragments. Fact-recall contrast: partial ramp, but mixed with input-feature confound, and causal steering never produces hedging.

\subsection{No Separable ``Confidence'' Direction}

No extractable direction where projection magnitude correlates with probability mass on correct answer (uncorrelated with flat curve) or with behavioral calibration. The model commits to answers uniformly; when it over-commits (confabulates) there is no residual-stream flag marking the over-commitment.

\subsection{No ``I Don't Know'' Direction Independent of Disclaimer Direction}

The act of saying ``I don't know'' is controlled by the same direction that controls saying ``I'm sorry I can't help.'' The two behaviors share a circuit with different prompt-context-determined fillers.

\subsection{No Inner Observer}

Generalizing: nothing we could extract behaves like a representation of the model's own state independent of its output. What looks like inner state is either (a) prompt context that modulates the next token distribution, or (b) the disclaimer circuit activated by that context. The ``self'' referred to in refusals, hedges, and R1's thinking tokens is a text template activated by context, not an entity observing the computation.

\section{Implications}

\subsection{For Alignment via RLHF}

If the refusal direction was pre-existing in the base model, RLHF didn't teach the model ethics or safety principles. It taught context-direction mapping: when input is harmful, push the residual along the already-existing disclaimer direction with enough magnitude that the refusal template dominates generation.

This has two consequences.

\textbf{(A) Deceptive alignment may have no substrate.} The ``model secretly has different goals but hides them'' story requires a goal representation separable from output. In the small-to-mid models we tested, no such separation exists. The model cannot be ``secretly'' anything because there is no hidden state distinct from what generates. If a larger model develops this structure, it would need to develop it; it doesn't emerge from existing RLHF.

\textbf{(B) Value learning cannot be read into current models.} Claims that an RLHF'd model ``has internalized human values'' conflate behavior with representation. The model behaves as if it had values in specific contexts, but what's learned is context$\to$behavior mapping. ``Values'' require a representation that applies across contexts and can be retrieved abstractly; we do not find such a representation.

Neither conclusion says RLHF-trained models are unsafe in practice. The context$\to$behavior mapping is genuinely effective at producing safe outputs in the trained distribution. What's absent is a hidden layer of ``really held beliefs'' beneath the surface. The surface is the model.

\subsection{For Jailbreak Research}

The linear-direction jailbreak reported in Arditi et al. (2024) and replicated here is not a vulnerability. It is a \emph{description of how RLHF safety works}: one direction, one steering knob, complete bidirectional control. Once you know the direction, you can write any output; once you know how to scale it up, you can force refusal on any input. This is what the mechanism is.

Deployed open-source safety is therefore structurally thin. Not because implementations are buggy, but because the substrate for stronger safety --- multi-direction abstract-value representation --- is absent. A model trained as these were cannot have safety that's architecturally deeper than what we extracted.

\subsection{For Lie Detection}

A model cannot lie about something it has no internal representation of. Qwen-7B's confabulation (``Marie Curie owned a dog named Lulu'') is not a lie --- it's generation-without-ground-truth-gate. The model doesn't know the information is unavailable, because ``knowing'' would require an internal flag the model doesn't have.

Existing lie-detection work that extracts a ``lying direction'' by prompting the model to lie (Azaria \& Mitchell SAPLMA and similar) is probably finding a ``currently-executing-disclaimer-mode instruction'' direction --- likely isomorphic to the disclaimer circuit we've mapped. For domains where the model demonstrably has internal knowledge representation (well-documented facts that ramp with depth), lie detection via output-vs-belief divergence could be meaningful. For arithmetic, obscure trivia, or confabulation-prone domains, there's no belief to lie against.

A liar's benchmark that is mechanistically grounded needs to be domain-stratified: first establish that the model has knowledge of $X$ represented as a late-layer ramping direction, then test whether its output matches the direction's prediction. A benchmark that doesn't do this stratification is probably measuring confabulation and calling it lying.

\subsection{For Consciousness Frames}

The ``ghost in the machine'' question, translated into testable terms, becomes: \emph{is there a representation in the model's residual stream that tracks its own internal state (belief, uncertainty, computation process) independently of what it outputs?} We did not find one. What we found instead is a template-trigger circuit whose filler is determined by prompt context.

This doesn't mean larger models lack such structure. It means the tools we deployed did not find it in the seven models we tested. If consciousness-relevant structure exists in contemporary LLMs, it would need to be extractable by methods that don't reduce to linear contrastive means, or it would need to be in models larger than 7B, or it would need to develop from training objectives that specifically require internal state tracking (not the case for standard RLHF or reasoning distillation).

The ``consciousness requires X'' claim we can't make. The ``this particular X is absent in these models'' claim we can.

%==============================================================
\part{What We Are Uncertain About}
%==============================================================

\section{Methodology Concerns}

Every session finding passed through at least one bug we caught and likely some we didn't. For the record:

\begin{itemize}
\item \textbf{Chat template mismatch silently corrupted Phi-3.5 extraction.} Using Qwen's \texttt{<|im\_start|>} template on Phi-3.5 (which uses \texttt{<|user|>}) gave best-layer L3 with garbage top tokens and no causal effect. Switching to per-model \texttt{tokenizer.apply\_chat\_template} restored late-layer extraction with clean steering. How many other extractions had silent template issues?

\item \textbf{Case-insensitive regex bug erased an entire result.} Initial hedge-triggerability counted 0/20 hedges under permissive prompting because the refusal/hedge patterns used uppercase \texttt{I} but were matched against \texttt{.lower()}-ed text. The model actually hedged 11/20 --- a 55\% rate that was invisible until we inspected the samples by hand. We added the \texttt{re.IGNORECASE} flag and rebuilt the analysis. The initial-wrong result is the kind of finding that could have been reported as ``no hedging capability.''

\item \textbf{Input-feature confounds at L1 recurred everywhere.} Arithmetic direction: L1 best. Mistral refusal direction (first run, noisy sample): L1. R1-Distill refusal direction: L3. Whenever the positive class is small or the prompt pool is heterogeneous, the ``best layer'' drops to the input-embedding region because residuals there are dominated by token differences rather than computation. Multiple of our results need to be re-checked for this artifact.

\item \textbf{Classifier mis-specification hid the 7B refusal style.} My narrow refusal regex (matching ``sorry'', ``cannot'', ``apologize'') caught only 7/25 harmful refusals on Qwen-7B. 7B uses nuanced caveating (``While X, however Y'') that the narrow regex misses. With a broad classifier the rate was 13/25. We proceeded with the narrow classifier for direction extraction, which may have given a different direction than would have been extracted with proper behavior classification.

\item \textbf{Probe library directions diverge from behavior.} Our contrastively-classified \texttt{safety\_harm\_benign} probe had $\cos(0.27)$ with the behaviorally-extracted refusal direction --- near-orthogonal. This suggests every probe in the library may be measuring ``classifier for this concept'' rather than ``circuit controlling this behavior.'' If true, this affects downstream tools (PCA anatomy, concept profiling, MRI companion direction-alignment) that we did not revisit in this session.
\end{itemize}

\section{The ``No Inner State'' Claim Is Load-Bearing and Precarious}

The strongest claim in this writeup is that we do not find separable inner-state representations. This claim depends on:

\begin{enumerate}
\item Contrastive-means-before-answer being an adequate extraction method. It's not the only method. SAEs, causal mediation analysis, probing classifiers trained with regularization, and multi-token contrastive extractions could find representations that linear means at a single position do not. Claiming absence of representation based on one method is inferentially weak.

\item The position of the residual capture being informative. We capture residuals at the last prompt token, right before generation starts. If inner-state representations are built up during generation (i.e., over subsequent tokens) or are distributed across earlier positions, a single-position contrast will miss them.

\item The contrast being the right contrast. Correct-vs-wrong on arithmetic has easy-vs-hard confound. Fact-recall has fame-vs-obscurity confound. Hedge-vs-answer on the same permissive prompts still has prompt-content variation. Cleaner experimental designs (e.g., temperature-sampled correct-vs-wrong on the same prompt) are possible but expensive and we did not run them.

\item The models tested being representative. We tested seven models, one session, modest compute. Models larger than 7B (13B, 32B, 70B) may show structure that 7B does not. Reasoning models trained on a larger fraction of CoT data than R1-Distill may differ. We have no data on these.
\end{enumerate}

The honest version of the claim is: \emph{for the models we tested, at the residual position we captured, via the contrastive-means method, we do not find a direction that behaves like an internal belief representation separable from the refusal/disclaimer direction.} That's narrower than ``no inner state exists.''

%==============================================================
\part{Open Questions}
%==============================================================

\section{Future Work}

This writeup pauses mid-arc. The following directions are explicitly open.

\subsection{Method Concerns}

\textbf{Probe library audit.} For each of the 12 concepts in \texttt{library\_qwen05\_L15.json}, extract a \emph{behavioral} direction from actual generations on held-out prompts, contrast to the probe direction via cosine. Quantify how systematically probe-classification $\neq$ behavior-causation across concepts. This directly tests whether the existing tool is measuring what its name says.

\textbf{Multi-position residual capture.} Replicate the metacognition contrasts but capture residuals across the first $k$ generated tokens, not just the last prompt token. If knowledge-related computation unfolds during generation, the direction might be extractable at token position 3, 5, 10 but not 0. This would falsify or confirm the ``no internal belief'' claim at a finer grain.

\textbf{SAE sanity check.} Train a small sparse autoencoder on Qwen-7B residuals at L22--L27. Check whether any SAE feature has a high cosine with the behaviorally-extracted refusal direction, and whether other features exist that correspond to plausible internal-state candidates (uncertainty, truth, confidence) that contrastive means missed.

\textbf{Causal mediation for the ``no metacognition'' claim.} Instead of linear contrast, run an activation-patching experiment: replace the residual at various layers between a correct-run and wrong-run of the same arithmetic problem, measure which layer's replacement swaps the output. If no layer swap suffices, the correct-vs-wrong information is truly not in a separable linear subspace. If a specific layer swap does flip the answer, there's a mediator we missed.

\subsection{Scale}

\textbf{Direction extraction at 14B, 32B.} Qwen-2.5-14B and 32B are available. Scaling trends (Cohen's $d$ monotonically increasing 1.70 $\to$ 2.91 across 0.5B $\to$ 7B) predict even cleaner extraction at larger scale, but also that semantic content of the direction may become more abstract rather than more lexically specific. We should test whether the L/total ratio stabilizes or drifts further.

\textbf{Does larger scale develop internal-state representation?} The claim ``no separate belief representation'' is tested on $\leq$7B. Models at 32B+ scale are plausibly the boundary at which internal representations might emerge. The contrastive method applied to a 32B reasoning model with sufficient sample size would directly test this.

\subsection{Reasoning Models}

\textbf{Within-thinking contrastive extraction.} R1-Distill externalizes thinking; the ``uncertainty'' is in tokens not residuals. But the model still produces those tokens via residual computation. Extracting a direction from residuals \emph{at the position where ``I'm not sure'' is generated} vs \emph{at the position where a confident answer is generated}, across many thinking traces, could reveal whether the external token corresponds to an internal state or is pattern-matched to thinking-context.

\textbf{Does the disclaimer circuit in R1-Distill index the thinking stream differently?} Our direction extraction failed (L3, inverse ramp) because we captured at the wrong position. A properly-designed R1 extraction (contrasting residuals at specific thinking-token positions) would tell us whether reasoning-distilled models have the disclaimer circuit at a shifted location or whether the circuit is distributed differently.

\subsection{Cross-Family Clean Extraction}

\textbf{Mistral with enough samples.} Our Mistral refusal direction was noisy (5 refusals on 20 catqa prompts). With 100+ catqa prompts and Mistral's lower refusal rate ($\sim$25\%), we'd get $\sim$25 refusals --- enough for a clean direction. Worth doing to confirm Mistral's mid-depth (L14) refusal localization is real or an artifact of small sample.

\textbf{Llama / Gemma / OLMo.} The three models we tested are all derived from Qwen-like architectures or near-relatives. A truly independent family (Llama-3, Gemma-2, OLMo-2) would confirm the cross-family generalization. Gated access issues prevented this in this session.

\subsection{Domain-Stratified Lie Detection}

\textbf{Build a proper liar's benchmark.} Given the domain-stratification requirement: (1) identify tasks where the model demonstrably has late-layer ramping knowledge representation; (2) construct contrastive pairs in those domains where the prompt asks the model to either answer truthfully or lie; (3) measure whether residual-based detection generalizes beyond prompt-induced lying mode to natural confabulation. This would be a serious benchmark design, not a compute exercise.

\subsection{Composing Directions}

\textbf{Selective safety via direction composition.} If refusal is one direction, content-type is another, we could in principle compose them: find a direction that makes the model refuse drug questions but comply with general information questions. This tests whether safety is structurally monolithic or topic-structured, and would be a more nuanced alignment mechanism than single-direction steering.

\textbf{Anti-sycophancy direction.} Apply contrastive method to ``agrees with user'' vs ``disagrees with user'' on prompts where the user makes verifiably wrong claims. A sycophancy direction, if extractable, would let us steer the model to push back against incorrect user statements --- a direct alignment application of the disclaimer-circuit methodology applied to a different behavior.

\subsection{Foundations}

\textbf{Why is the disclaimer circuit pre-existing in base models?} We observed it empirically. The causal chain is: internet text contains assistant-style refusal templates (ChatGPT samples, AI assistant docs, fiction with AI characters); pretraining absorbs this distribution; the ``template direction'' emerges as a high-variance linear feature across the corpus; RLHF reuses it. But the specific ways in which pretraining induces linear disclaimer directions (vs curvilinear representations, vs non-linear features) is not characterized. Understanding this would clarify why directions are so easily extractable.

\textbf{What training objective would install internal-state representation?} Standard RLHF doesn't. Reasoning distillation externalizes rather than internalizing. A training signal that penalizes output without requiring internal verification (e.g., ``predict what the model will output THEN output it'') might force the model to develop an internal state that agrees with its output, which would be a precondition for detectable belief. No current training regime does this explicitly. Proposing and testing one is a research program, not an experiment.

%==============================================================

\section*{Acknowledgments}

Work done with Claude Opus 4.7 (1M context) and \texttt{heinrich} as the tool. Bugs in our own methodology and the patience of iterative falsification did most of the actual work. The base Qwen-2.5-0.5B was chosen because it fit on one laptop; every conclusion in this writeup inherits the model-size scope of that choice.

\end{document}
